\emph{Setting.} $\Cl(2,4)$ is realised on orthonormal generators $e_{0}, \dots, e_{5}$ with squares $(+,-,-,-,+,-)$, $e_{0}, \dots, e_{3}$ the spacetime vectors in the store signature, $\nO = \tfrac12(e_{4} + e_{5})$, $\nInf = e_{4} - e_{5}$ (the chart of the stored check of Theorem \claimref{dark-energy.frozen-point-makes-stelle-west-vector}); then $\nO \cdot \nInf = 1$, both are null, $\nInf \wedge \nO = e_{4}e_{5}$ and $i_{6} = e_{0}e_{1}e_{2}e_{3}e_{4}e_{5}$. In the chart with $\nO$ and $\nInf$ built from $e_{4} - e_{5}$ and $e_{4} + e_{5}$ the other way round, with the same $i_{6}$, every $\epsilon$ changes sign. For a bivector $B$ and a vector $a$, $B \cdot a = \tfrac12(Ba - aB)$. A bivector-valued two-form is, in Gauge Theory Gravity, a linear map $B \mapsto F(B)$ from physical bivectors to bivectors of $\Cl(2,4)$, with frame components $F_{ef} = F(e_{e}\wedge e_{f})$ (the field strength of $\Omega_{T}$ on the physical planes, Theorem \claimref{synthesis.gtg-tractor-gauge-field-reproduces-tractor-connection}); the four-form of a bilinear pairing $\beta(F \wedge G)$ has frame coefficient $\tfrac14\sum_{p}\sgn(p)\,\beta(F_{p_{0}p_{1}}, G_{p_{2}p_{3}})$ over the permutations $p$ of $(0,1,2,3)$, equivalently the sum over ordered pairs of complementary coordinate planes (the convention of $F = \tfrac12 F_{ef}\,e^{e} \wedge e^{f}$). Since the statements (i) to (iv) are pointwise and algebraic, it suffices to work at a point in an orthonormal frame of the derived metric.

\emph{(i) The locus.} By Theorem \claimref{synthesis.gtg-tractor-gauge-field-reproduces-tractor-connection}(iv) the field strength of $\Omega_{T}$ has spacetime block $-\mathcal{W}(A \wedge B)$ (the algebraic part) and two derivative blocks: the torsion along $\nInf$, zero for a torsion-free field, and the Cotton term along $\nO$. Each value is therefore $-\mathcal{W}(A \wedge B) + \nO \wedge c$ with $c$ a spacetime vector. A spacetime bivector is orthogonal to $\nO$, so $\mathcal{W}(A \wedge B)\cdot\nO = 0$, and $(\nO \wedge c)\cdot\nO = \nO(c \cdot \nO) - c\,(\nO \cdot \nO) = 0$. A torsion value $\nInf \wedge t$ gives $(\nInf \wedge t)\cdot\nO = -t \neq 0$. The span of the values is contained in $\Lambda^{2}(X^{\perp})$, $X^{\perp} = \mathrm{span}(e_{0}, \dots, e_{3}, \nO)$, a five-dimensional degenerate subspace containing $X$. The check builds the field strength from three sample Weyl tensors (Weyl parts, $R - \mathrm{KN}(\Sch, g)$ with the Schouten tensor of Fact \claimref{synthesis.imported-gtg-dictionary}, of Kulkarni-Nomizu sums of fixed rational symmetric matrices, verified Ricci-flat and nonzero) and a fixed rational Cotton block, and verifies $FX = XF$ on every frame component, and the failure with a torsion block.

\emph{(ii) Independence of $Y$.} For $F, G \in \Lambda^{2}(X^{\perp})$, $X \wedge F \wedge G \in \Lambda^{5}(X^{\perp})$, which is one-dimensional. If $Z \in X^{\perp}$ then $X \wedge F \wedge G \wedge Z \in \Lambda^{6}(X^{\perp}) = 0$. Writing any $Y$ with $\ip{X}{Y} \neq 0$ as $Y = \ip{X}{Y}\nInf + Z$ with $Z \in X^{\perp}$ (possible because $\nInf \notin X^{\perp}$ and $\ip{X}{\nInf} = 1$) gives $\epsilon(Y, X, F, G) = \ip{X}{Y}\,\epsilon(\nInf, X, F, G)$, so $\epsilon_{X}(F, G) := \epsilon(Y, X, F, G)/\ip{X}{Y}$ is well defined at fixed $X$; it is homogeneous of degree zero in $X$ and is not a polynomial in $(X, F)$ (an $\so(2,4)$-invariant polynomial of degree zero in $X$ is a function of $F$ alone). By antisymmetry of the wedge, $\epsilon(I, V, F, G) = \epsilon(I, X, F, G) - \tfrac{3}{\Lambda}\epsilon(I, I, F, G) = \epsilon(I, X, F, G) = \ip{X}{I}\,\epsilon_{X}(F, G)$; the blade identity $I \wedge V = I \wedge X$ is that of the stored lemma (also \claimref{synthesis.frozen-point-action-is-bf-transport-broken-by-the-blade-i-wedge-x}(i)). The check confirms the independence with three choices of $Y$ (including components along $\nO$ and along spacetime vectors), the identity with $I = \nInf + \tfrac{\Lambda}{6}\nO$ at $\Lambda \in \{3, -3, \tfrac97\}$ (with $Q(I) = \Lambda/3$, $Q(V) = -3/\Lambda$, $\ip{I}{V} = 0$ as in Theorem \claimref{dark-energy.frozen-point-makes-stelle-west-vector}) and with a general tractor $\tfrac52\nInf + e_{1} - 4\nO$ (where the value is $\tfrac52\epsilon_{X}$), and that with a torsion block two choices of $Y$ give different values.

\emph{(iii) The Weyl square.} Only the Lorentz block contributes: a term $\nO \wedge c$ in either slot gives $X \wedge \nO = 0$. On the Lorentz block, $\nInf \wedge X = e_{4}e_{5}$ and $\epsilon_{X}(F, G)$ is the coefficient of $e_{0}e_{1}e_{2}e_{3}$ in $F \wedge G$ (the four-dimensional fibre volume form). Both sides of $\epsilon_{X}(F, F) = \tfrac14 W_{abcd}W^{abcd}$ are quadratic in the Weyl tensor, so it suffices to verify the polarised identity $\epsilon_{X}(F_{1}, F_{2}) = \tfrac14 (W_{1})_{abcd}(W_{2})^{abcd}$ on a basis of the Weyl space. In Gauge Theory Gravity the Weyl part is the linear map $\mathcal{W}$ on physical bivectors, the Lorentz block is $F(B) = -\mathcal{W}(B)$, and the identity reads $\epsilon_{X}(F_{1}, F_{2}) = \tfrac12\langle\mathcal{W}_{1}(a\wedge b)\,\mathcal{W}_{2}(\partial_{b}\wedge\partial_{a})\rangle_{0}$, a bilinear form on the ten-dimensional space of Weyl maps (trace-free, $\partial_{a}\cdot\mathcal{W}(a\wedge b) = 0$, and $I$-linear). The check builds ten independent Weyl maps with the library, as the Weyl parts of the Kulkarni-Nomizu operators $a\wedge b \mapsto S(a)\wedge T(b) + T(a)\wedge S(b)$ of pairs of basis symmetric maps, confirms that they are trace-free and $I$-linear, and verifies the polarised identity on all fifty-five pairs and with the Cotton block present (check\_gtg\_weyl\_square\_of\_weyl\_maps); in components the same identity is $\epsilon_{X}(F_{1}, F_{2}) = \tfrac14(W_{1})_{abcd}(W_{2})^{abcd}$, which the earlier check verifies on a basis of ten component Weyl tensors (rank ten, the dimension of the Weyl space in four dimensions) and all fifty-five pairs; it also verifies it on the samples with the Cotton block present, and that it fails on a curvature with a Ricci part (so the identity is specific to the Weyl block). This is the classical fact that the internal-epsilon contraction of the Weyl curvature two-forms is the Weyl square, the Weyl tensor being equal to minus its double dual in Lorentzian signature; read on the conformal-group curvature with the conventional constraints it is the construction of conformal gravity by Kaku, Townsend and van Nieuwenhuizen 1977 (Fact \claimref{synthesis.imported-conformal-gravity-bach-equations-and-conformal-gauge-form}(b)), and the identity is theirs. The series-specific content is (ii), the independence of the second vector on the locus, which makes the frozen point drop out. With the prefactor $1/Q(I) = 3/\Lambda$ of Theorem \claimref{dark-energy.frozen-point-makes-stelle-west-vector}(iii) and (ii), the functional on transport is $\tfrac{3}{4\Lambda}\ip{X}{I}\int W_{abcd}W^{abcd}\,\mathrm{vol}$, and $\ip{X}{I} = 1$ for the frozen point; divided by $16\pi G$ and with $\ell^{2} = 3/\Lambda$ this is $(\ell^{2}/64\pi G)\int W_{abcd}W^{abcd}\,\mathrm{vol}$, the value recorded in Theorem \claimref{synthesis.promoted-norm-weyl-clock-flat-weight} for Einstein instantons, here for every torsion-free field. On Einstein fields this agrees with Fact \claimref{synthesis.imported-miskovic-olea-identity} (there the shifted curvature is the Weyl part); the present statement holds for every torsion-free field.

\emph{(iv) The quadratic densities on the locus.} The stabiliser of $X$ in the bivectors is $\Lambda^{2}(X^{\perp})$ itself (ten-dimensional), since a bivector commutes with $X$ exactly when it annihilates it. Invariance of a bilinear form $\beta$ under it is the linear system $\beta([\xi, u], v) + \beta(u, [\xi, v]) = 0$, $[\xi, u] = \tfrac12(\xi u - u\xi)$, for $\xi$ and $u, v$ in a basis of $\Lambda^{2}(X^{\perp})$. The check solves it exactly: the symmetric solutions form a two-dimensional space and the antisymmetric solutions are zero. $K$ and $\epsilon_{X}$ satisfy the system (checked) and are linearly independent (checked), so they span it. Under the internal reflection $x \mapsto e_{1}\hat{x}e_{1}^{-1}$, which fixes $X$, with the spacetime orientation fixed, $K$ is unchanged and $\epsilon_{X}$ changes sign (checked on a pair with $\epsilon_{X} \neq 0$). Parity in the usual sense reverses the spacetime orientation as well, which reverses the sign of every four-form integral; so $\epsilon_{X}(F, F)$, the Weyl square, is parity-even and $K(F \wedge F)$, on the locus the Pontryagin density of the Weyl curvature, is parity-odd. For contrast, the same system on all fifteen bivectors under the whole algebra has a one-dimensional symmetric solution space (the Killing form), and on the Lorentz bivectors of $\Cl(1,3)$ the scalar and pseudoscalar parts of the product are two independent invariant forms (Pontryagin and Euler). The Chern form is exact: for any connection $A$ with $F = dA + A \times A$ in frame components ($F_{\mu\nu} = \partial_{\mu}A_{\nu} - \partial_{\nu}A_{\mu} + A_{\mu} \times A_{\nu}$), $\langle F \wedge F\rangle_{0} = d\,\mathrm{CS}(A)$ with $\mathrm{CS}(A) = \langle A \wedge dA\rangle_{0} + \tfrac13\langle A \wedge (A \times A)\rangle_{0}$, by the Bianchi identity (Chern and Simons 1974, Fact \claimref{synthesis.imported-chern-weil-and-chern-simons-forms}(b)); this is a four-dimensional statement, since only quadratic invariants fill four form degrees (in five dimensions the cubic Pfaffian has a Chern-Simons form with field equations, Theorem \claimref{synthesis.macdowell-mansouri-is-the-frozen-point-reduction-of-five-dimensional-chern-simons}); the check verifies the identity exactly at a point on a fixed second-order jet of a $\Cl(2,4)$ connection, and that the coefficients $\tfrac23$ and $0$ fail. The frame legs: $DX$ takes values in $X^{\perp}$ ($\ip{X}{DX} = \tfrac12 d\,Q(X) = 0$), and in the normal gauge $DX$ is the frame, so $DX \wedge DX$ has values $2\,e_{e} \wedge e_{f}$. Then $K(DX \wedge DX, F)$ is the contraction $e^{a} \wedge e^{b} \wedge W_{ab}$, zero by the cyclic identity of the Weyl tensor, and $\epsilon_{X}(DX \wedge DX, F)$ is the Einstein-Hilbert-type contraction, proportional to the scalar curvature of the Weyl block, zero. Both pairings are linear in $W$; the check verifies both on all ten basis Weyl tensors and on the samples (with Cotton block), which is complete, and that $\epsilon_{X}(DX \wedge DX, F)$ is nonzero for a curvature with a Ricci part and $K(DX \wedge DX, F)$ nonzero for a curvature two-form violating the cyclic identity. The remaining frame-leg density $\epsilon_{X}(DX \wedge DX, DX \wedge DX)$ is the volume form (nonzero, checked); it has conformal weight four, so it is not a functional of the conformal class that the normal tractor connection carries (it needs a scale, which the frozen point supplies).

\emph{(v) Not topological.} If a density $L$ had identically vanishing Euler-Lagrange derivative, then on a closed manifold $\int L$ would be constant along every smooth family of metrics. On $S^{2}(K_{1}) \times S^{2}(K_{2})$ (Riemannian, curvature $K_{1}$ on the first factor's planes and $K_{2}$ on the second's) the curvature operator at each point is, on the bivectors of $\Cl(4,0)$, $\mathcal{R}(e_{1}\wedge e_{2}) = K_{1}\,e_{1}\wedge e_{2}$, $\mathcal{R}(e_{3}\wedge e_{4}) = K_{2}\,e_{3}\wedge e_{4}$ and zero on the four mixed planes; its Weyl map and the Euler density $2\langle\mathcal{R}(a\wedge b)\mathcal{R}(\partial_{b}\wedge\partial_{a})\rangle_{0} - 4\mathcal{R}(\partial_{a})\cdot\mathcal{R}(a) + \mathcal{R}^{2}$ are computed pointwise, $2\langle\mathcal{W}(a\wedge b)\mathcal{W}(\partial_{b}\wedge\partial_{a})\rangle_{0} = \tfrac43(K_{1} + K_{2})^{2}$ and $8K_{1}K_{2}$ (check\_gtg\_weyl\_square\_of\_weyl\_maps; the earlier check computes the same from component tensors). Beyond one flat chart: the closed product manifold and the integrals over it are global and are not carried by a single gravitational field on flat spacetime; the pointwise densities are, and the integrals are read on the manifold assembled from the derived metric's charts. Integrated: per $16\pi^{2}$ (the product of the areas $4\pi/K_{i}$ is $16\pi^{2}/(K_{1}K_{2})$), $\int W_{abcd}W^{abcd} = \tfrac43(K_{1} + K_{2})^{2}/(K_{1}K_{2})$, equal to $\tfrac{16}{3}$ at $K_{1} = K_{2}$ and $6$ at $K_{2} = 2K_{1}$, while $\int(R_{abcd}R^{abcd} - 4R_{ab}R^{ab} + R^{2}) = 8$ for both (Gauss-Bonnet: $32\pi^{2}\chi$ with $\chi = 4$). The family is smooth in $K_{2}/K_{1}$, so the Euler-Lagrange derivative of the Weyl square is not identically zero. That derivative is a universal polynomial in the metric, its inverse and their derivatives (times a power of $|\det g|$); if it vanished on all Lorentzian metrics it would vanish on an open set of jets and hence identically, also in Riemannian signature, which the computation excludes. So the Weyl-square density is not topological in the store's signature either; its Euler-Lagrange derivative is the Bach tensor (Bach 1921, Fact \claimref{synthesis.imported-conformal-gravity-bach-equations-and-conformal-gauge-form}(a)), not used in this argument.

\emph{Reading.} By (i) to (iv), with the canonical tractor alone and the field strength on transport's locus, every invariant density quadratic in $F$ is $\alpha K + \beta\epsilon_{X}$, $K$ exact in four dimensions, and the frame-leg pairings vanish, so no density in this class has an Einstein-Hilbert term there; by (v) $\epsilon_{X}$ is a genuine action, conformal gravity. By (ii) and (iii) the stored frozen-point functional, evaluated on transport, is this action with coupling $\tfrac{3}{4\Lambda}$; the frozen point contributes the prefactor and nothing else. The Einstein equations of Theorem \claimref{synthesis.macdowell-mansouri-principle-fixes-the-trace-free-part} arise because that functional is varied over connections preserving $I$ (Theorem \claimref{dark-energy.frozen-point-makes-stelle-west-vector}(iii)), which leave the locus (they carry torsion and a Schouten slot not tied to the curvature), and for the normal tractor connection preserving $I$ is the Einstein condition (Theorem \claimref{synthesis.gtg-tractor-gauge-field-reproduces-tractor-connection}(ii)). The hypothesis that admits only the $I$-preserving connections is the absence of an inverse frame, which excludes the Schouten slot of the normal connection (Theorem \claimref{synthesis.curvature-only-tractor-actions-are-macdowell-mansouri}(vi), cited and not used as a premise). Restricted to a conformal class, by (ii) the integrand does not see the scale beyond $\ip{X}{I}$, which extends Theorem \claimref{synthesis.positive-structure-and-natural-scale-actions-do-not-select-the-almost-einstein-scale} to this integrand. The class is limited: with a scale tractor treated as a field and its derivative admitted, tractor actions with an Einstein-Hilbert term are known (cited from memory); they are outside the class of (iv). That Einstein fields solve the Bach equations, and that conformally Einstein fields do too, is Fact \claimref{synthesis.imported-conformal-gravity-bach-equations-and-conformal-gauge-form}(a); it is used only in this Reading, for the statement that the integrand's field equations on transport are the Bach equations and not Einstein's.
